\clearpage
\hypertarget{native-264-pure13}{}\label{native-264-pure13}
\Article{13}
Now that we have shown how to expand a product of two factors, expanding a product of more factors will present no difficulty. Once the product $[r,R^\mu]\times[r,R^\nu]$ has been reduced to the form $A[1,R^{\mu+\nu}]+B[r,R^{\mu+\nu}]$, it is evident that, on introducing a third factor $[r,R^\pi]$, the product becomes $=C[1,R^{\mu+\nu+\pi}]+D[r,R^{\mu+\nu+\pi}]$ if we put
\[
[r,R^{\mu+\nu}][r,R^\pi]=c[1,R^{\mu+\nu+\pi}]+d[r,R^{\mu+\nu+\pi}]
\]
and
\[
\begin{aligned}
C&=Bc\\
D&=Bd+A\{1+R^{\mu+\nu}+R^{2\mu+2\nu}+\text{etc.}+R^{(n-2)(\mu+\nu)}\}
\end{aligned}
\]
Thus $[r,R]^\lambda$ can easily be reduced to the form $A[1,R^\lambda]+B[r,R^\lambda]$.

As an example, we shall expand the powers of the function $[r,R]$ for $n=11$, $\beta=5$, setting $g=2$. The following indices then correspond to the numbers:
\[
\begin{array}{r@{\quad}r@{\pdot}r@{\pdot}r@{\pdot}r@{\pdot}r@{\pdot}r@{\pdot}r@{\pdot}r@{\pdot}r@{\pdot}r}
\text{numbers}&1&2&3&4&5&6&7&8&9&10\\
\text{indices}&0&1&8&2&4&9&7&3&6&5
\end{array}
\]
Thus, to expand the square $[r,R]^2$ by formula I of Article~11, we have
\[
\mu=1,\quad\nu=1
\]
\[
\begin{array}{l@{\quad}r@{\pdot}r@{\pdot}r@{\pdot}r@{\pdot}r@{\pdot}r@{\pdot}r@{\pdot}r@{\pdot}r}
\text{values of }x\ \ldots&1&2&3&4&5&6&7&8&9\\
\operatorname{ind}x\ \ldots&0&1&8&2&4&9&7&3&6\\
2\operatorname{ind}(x+1)\ \ldots&2&16&4&8&18&14&6&12&10\\
\multicolumn{10}{l}{\text{Least residues of }\operatorname{ind}x-2\operatorname{ind}(x+1)}\\
\quad\text{modulo }5\ \ldots&3&0&4&4&1&0&1&1&1
\end{array}
\]
from which we deduce
\SourceLeaf{265}
\[
\Omega=[r,R^2]\times\{2+4R+R^3+2R^4\}
\]
and
\[
1^0.\qquad[r,R]^2=[1,R^2]+[r,R^2]\times\{2+4R+R^3+2R^4\}
\]
The same expression follows from formula III of Article~11, namely
\[
\begin{array}{l@{\quad}r@{\pdot}r@{\pdot}r@{\pdot}r@{\pdot}r@{\pdot}r@{\pdot}r@{\pdot}r@{\pdot}r}
\text{values of }z\ \ldots&2&3&4&5&6&7&8&9&10\\
\operatorname{ind}z\ \ldots&1&8&2&4&9&7&3&6&5\\
\operatorname{ind}(n+1-z)\ \ldots&5&6&3&7&9&4&2&8&1\\
\multicolumn{10}{l}{\text{least residues of }\operatorname{ind}z+\operatorname{ind}(n+1-z)}\\
\quad\text{modulo }5\ \ldots&1&4&0&1&3&1&0&4&1
\end{array}
\]

In exactly the same way we find
\[
\begin{aligned}
2^0.\qquad[r,R^2]\pdot[r,R]&=[1,R^3]+[r,R^3]\times\{2+R+4R^2+2R^3\}\\
3^0.\qquad[r,R^3]\pdot[r,R]&=[1,R^4]+[r,R^4]\times\{2+4R+R^3+2R^4\}
\end{aligned}
\]
Finally, we obtain
\[
4^0.\qquad[r,R^4]\pdot[r,R]=[1,1]+[r,1]\times\{1+2R+2R^2+2R^3+2R^4\}
\]
Multiplying equation $1^0$ by $[r,R]$ and substituting for $[r,R^2]\pdot[r,R]$ its value from $2^0$, and using
\[
[1,R^2]\pdot[r,R]=[r,R^3]\pdot\{2+2R+2R^2+2R^3+2R^4\}
\]
we deduce
\[
\begin{aligned}
\relax[r,R]^3&=[1,R^3]\times\{2+4R+R^3+2R^4\}\\
&\quad+[r,R^3]\times\{12+22R+18R^2+24R^3+15R^4\}
\end{aligned}
\]
and in the same way
\[
\begin{aligned}
\relax[r,R]^4&=[1,R^4]\times\{12+22R+18R^2+24R^3+15R^4\}\\
&\quad+[r,R^4]\times\{164+170R+205R^2+180R^3+190R^4\}\\
\relax[r,R]^5&=[1,1]\times\{164+170R+205R^2+180R^3+190R^4\}\\
&\quad+[r,1]\times\{1836+1830R+1795R^2+1820R^3+1810R^4\}\\
&=1640+1700R+2050R^2+1800R^3+1900R^4\\
&\quad+(1836+1830R+1795R^2+1820R^3+1810R^4)(s-1)\\
&=918Ss-98S-(6R+41R^2+16R^3+26R^4)s\\
&\quad+66R+451R^2+176R^3+286R^4
\end{aligned}
\]

\SourceLeaf{266}
\Needspace{25\baselineskip}
\Article{14}
The calculation in the preceding articles\textsuperscript{'} was carried out so that it could be extended to all values of $r$ and $R$. It becomes considerably shorter if we regard $R$ from the outset as a proper root of the equation $x^\beta-1=0$. Under this assumption the product $[r,R^\mu]\times[r,R^\nu]$ reduces to the form $B[r,R^{\mu+\nu}]$ whenever $\mu+\nu$ is not divisible by $\beta$; but when $\mu+\nu$ is divisible by $\beta$, that product becomes $=(n-1)R^{\frac12(n-1)\mu}+[r,1]\Sigma R^{\mu\operatorname{ind}x}$, substituting for $\operatorname{ind}x$ all the numbers $0,1,2,3\dotsfour n-2$ except $\frac12(n-1)$. It follows easily that in this case, if $\mu$ and therefore also $\nu$ are not divisible by $\beta$, we have
\[
[r,R^\mu]\pdot[r,R^\nu]=R^{\frac12(n-1)\mu}\{n-1-[r,1]\}
\]
and hence $=0$ for $r=1$, and $=nR^{\frac12(n-1)\mu}$ for every other value of $r$. Moreover, since $R^{\frac12(n-1)\mu}$ becomes $=+1$ or $=-1$ according as $\frac{n-1}\beta\pdot\mu$ is even or odd, our product becomes $=n$ in the former case and $=-n$ in the latter.

It further follows that we can put
\[
\begin{aligned}
\relax[r,R]^2&=A'[r,R^2]\\
\relax[r,R^2]\pdot[r,R]&=A''[r,R^3]\\
\relax[r,R^3]\pdot[r,R]&=A'''[r,R^4]
\end{aligned}
\]
and so on, up to
\[
[r,R^{\beta-2}]\pdot[r,R]=A^{(\beta-2)}[r,R^{\beta-1}]
\]
from which we have
\[
\begin{aligned}
\relax[r,R]^2&=A'[r,R^2]\\
\relax[r,R]^3&=A'A''[r,R^3]\\
\relax[r,R]^4&=A'A''A'''[r,R^4]
\end{aligned}
\]
and so on. Finally,
\[
[r,R]^\beta=\pm nA'A''A'''\dots A^{(\beta-2)}
\]
where the upper or lower sign is to be taken according as $\frac{n-1}\beta$ is even or odd.

It is therefore evident that, once the value of $[r,R]$ has been found, the remaining functions
\[
[r,R^2]=\frac{[r,R]^2}{A'},\quad[r,R^3]=\frac{[r,R]^3}{A'A''}\quad\text{etc.}
\]
can be determined here much more readily than in the cases where the condition $\alpha$ $=1$ does not hold,
\SourceLeaf{267}
as we already noted in the \emph{Disq. Ar.} (Article~360, \textsc{iii}). A fuller consideration of the nature of the functions $A',A''$, etc., makes these operations still easier.

\Article{15}
In Article~9 we showed that the value of the function $[r,R]$ can be reduced to the form $\surd n(\cos f+i\sin f)$, and in the same way the functions $[r,R^2]$, $[r,R^3]$, etc., up to $[r,R^{\beta-1}]$, can be reduced to a similar form. Let us put
\[
\begin{aligned}
\relax[r,R]&=\surd n(\cos f'+i\sin f')\\
\relax[r,R^2]&=\surd n(\cos f''+i\sin f'')\\
\relax[r,R^3]&=\surd n(\cos f'''+i\sin f''')\\
&\text{etc.}
\end{aligned}
\]
Then
\[
\begin{aligned}
A'&=\surd n(\cos(2f'-f'')+i\sin(2f'-f''))\\
A''&=\surd n(\cos(f'+f''-f''')+i\sin(f'+f''-f'''))\\
A'''&=\surd n(\cos(f'+f'''-f'''')+i\sin(f'+f'''-f''''))\\
&\text{etc.}
\end{aligned}
\]
It is evident from this that, if the functions $A',A'',A'''$, etc., are reduced to the forms
\[
\begin{aligned}
A'&=a'(\cos b'+i\sin b')\\
A''&=a''(\cos b''+i\sin b'')\\
A'''&=a'''(\cos b'''+i\sin b''')\\
&\text{etc.}
\end{aligned}
\]
with all of $a',a'',a'''$, etc., positive, we shall have
\[
\begin{aligned}
a'&=a''=a'''\text{ etc.}=\surd n\\
f'&=\frac1\beta(b'+b''+b'''+\text{etc.}+b^{(\beta-2)})
\end{aligned}
\]
if $\frac{n-1}\beta$ is even, or
\[
f'=\frac1\beta(180^0+b'+b''+\text{etc.}+b^{(\beta-2)})
\]
if $\frac{n-1}\beta$ is odd, and then
\[
\begin{aligned}
\relax[r,R]&=\surd n(\cos f'+i\sin f')\\
\relax[r,R^2]&=\surd n(\cos(2f'-b')+i\sin(2f'-b'))\\
\relax[r,R^3]&=\surd n(\cos(3f'-b'-b'')+i\sin(3f'-b'-b''))\\
&\text{etc.}
\end{aligned}
\]
\SourceLeaf{268}
Finally, by the formulas of Article~8,
\[
\begin{aligned}
\left(\frac{n-1}\beta,1\right)&=-\frac1\beta+\frac{\surd n}\beta\{\cos f''+\cos(2f'-b')+\cos(3f'-b'-b'')+\text{etc.}\\
&\qquad+\cos((\beta-1)f'-b'-b''-b'''-\text{etc.}-b^{(\beta-2)})\}\\
&\quad+\frac{i\surd n}\beta\{\sin f'+\sin(2f'-b')+\sin(3f'-b'-b'')+\text{etc.}\}
\end{aligned}
\]
and similarly we obtain the values of the functions $\left(\frac{n-1}\beta,g\right)$, $\left(\frac{n-1}\beta,g^2\right)$, $\left(\frac{n-1}\beta,g^3\right)$, etc., if in this formula $f'$ is replaced respectively by $f'-\frac{360^0k}\beta$, $f'-2\frac{360^0k}\beta$, $f'-3\frac{360^0k}\beta$, etc., assuming $R=\cos\frac{360^0k}\beta+i\sin\frac{360^0k}\beta$.

\clearpage
\Article{16}
A further simplification follows from the observation below. Since by Article~14 we have
\[
\pm[r,R][r,R^{\beta-1}]=[r,R^2][r,R^{\beta-2}]=\pm[r,R^3][r,R^{\beta-3}]\text{ etc.}=n
\]
taking [in the first, third, etc., product] the upper or lower sign according as $\frac{n-1}\beta$ is even or odd, in the former case we must have
\[
\cos(f'+f^{(\beta-1)})=\cos(f''+f^{(\beta-2)})=\cos(f'''+f^{(\beta-3)})\text{ etc.}=1
\]
in the latter
\[
-\cos(f'+f^{(\beta-1)})=\cos(f''+f^{(\beta-2)})=-\cos(f'''+f^{(\beta-3)})\text{ etc.}=1
\]
and in both cases
\[
\sin(f'+f^{(\beta-1)})=\sin(f''+f^{(\beta-2)})=\sin(f'''+f^{(\beta-3)})\text{ etc.}=0
\]
We may therefore put, in the former case,
\[
f^{(\beta-1)}=-f',\quad f^{(\beta-2)}=-f'',\quad f^{(\beta-3)}=-f'''\quad\text{etc.}
\]
in the latter
\[
f^{(\beta-1)}=180^0-f',\quad f^{(\beta-2)}=-f'',\quad f^{(\beta-3)}=180^0-f'''\quad\text{etc.}
\]
From this it follows that, in the former case,
\[
\begin{aligned}
b^{(\beta-2)}&=b',&b^{(\beta-3)}&=b'',&b^{(\beta-4)}&=b'''\quad\text{etc.}\\
A^{(\beta-2)}&=A',&A^{(\beta-3)}&=A'',&A^{(\beta-4)}&=A'''\quad\text{etc.}
\end{aligned}
\]
\Needspace{21\baselineskip}
and in the latter
\SourceLeaf{269}
\[
\begin{aligned}
b^{(\beta-2)}&=b'-180^0,&b^{(\beta-3)}&=b''+180^0,&b^{(\beta-4)}&=b'''-180^0\quad\text{etc.}\\
A^{(\beta-2)}&=-A',&A^{(\beta-3)}&=-A'',&A^{(\beta-4)}&=-A'''\quad\text{etc.}
\end{aligned}
\]
so that the number of functions $A',A'',A'''$, etc., is reduced by half. It further follows that, in the former case,
\begingroup\postdisplaypenalty=10000
\[
f'=\frac1\beta(2b'+2b''+\text{etc.}+2b^{(\frac12\beta-1)})
\]
\endgroup
\begingroup\postdisplaypenalty=10000
\[
\begin{aligned}
\left(\frac{n-1}\beta,1\right)&=-\frac1\beta+\frac{\surd n}\beta\{2\cos f'+2\cos(2f'-b')\\
&\qquad+2\cos(3f'-b'-b'')+\text{etc.}\\
&\qquad+2\cos((\tfrac12\beta-1)f'-b'-b''-\text{etc.}-b^{(\frac18\beta-2)})\\
&\qquad+\cos(\tfrac12\beta f'-b'-b''-\text{etc.}-b^{(\frac12\beta-1)})\}
\end{aligned}
\]
\endgroup
(where the last term is evidently $=\cos0=1$) or
\begingroup\postdisplaypenalty=10000
\[
f'=\frac1\beta(2b'+2b''+\text{etc.}+2b^{(\frac12(\beta-3))}+b^{(\frac12(\beta-1))})
\]
\endgroup
\[
\begin{aligned}
\left(\frac{n-1}\beta,1\right)&=-\frac1\beta+\frac{\surd n}\beta\{2\cos f'+2\cos(2f'-b')\\
&\qquad+2\cos(3f'-b'-b'')+\text{etc.}\\
&\qquad+2\cos(\tfrac12(\beta-1)f'-b'-b''-\text{etc.}-b^{(\frac12(\beta-3))})\}
\end{aligned}
\]
according as $\beta$ is even or odd; and in the latter case,

\begingroup\postdisplaypenalty=10000
\[
f'=\frac1\beta(2b'+2b''+\text{etc.}+2b^{(\frac12\beta-1)})
\]
\endgroup
\[
\begin{aligned}
\left(\frac{n-1}\beta,1\right)&=-\frac1\beta+\frac{\surd n}\beta\{2\cos(2f'-b')+2\cos(4f'-b'-b''-b''')+\text{etc.}\\
&\qquad+2\cos((\tfrac12\beta-2)f'-b'-b''-\text{etc.}-b^{(\frac12\beta-3)})\\
&\qquad+\cos(\tfrac12\beta f'-b'-b''-\text{etc.}-b^{(\frac12\beta-1)})\}\\
&\quad+i\frac{\surd n}\beta\{2\sin f'+2\sin(3f'-b'-b'')+\text{etc.}\\
&\qquad+2\sin((\tfrac12\beta-1)f'-b'-b''-\text{etc.}-b^{(\frac12\beta-2)})\}
\end{aligned}
\]
\Needspace{32\baselineskip}
or
\[
f'=\frac1\beta(2b'+2b''+\text{etc.}+2b^{(\frac12\beta-1)}+180^0)
\]
\[
\begin{aligned}
\left(\frac{n-1}\beta,1\right)&=-\frac1\beta+\frac{\surd n}\beta\{2\cos(2f'-b')+2\cos(4f'-b'-b''-b''')+\text{etc.}\\
&\qquad+2\cos((\tfrac12\beta-1)f'-b'-b''-\text{etc.}-b^{(\frac12\beta-2)})\}\\
&\quad+i\frac{\surd n}\beta\{2\sin f'+2\sin(3f'-b'-b'')+\text{etc.}\\
&\qquad+2\sin((\tfrac12\beta-2)f'-b'-b''-\text{etc.}-b^{(\frac12\beta-3)})\\
&\qquad+\sin(\tfrac12\beta f'-b'-b''-\text{etc.}-b^{(\frac12\beta-1)})\}
\end{aligned}
\]
according as $\frac12\beta$ is even or odd. What we noted above also holds for the remaining periods of $\frac{n-1}\beta$ terms. In general, therefore, we conclude that determining these periods requires the division of the whole circle into $\beta$ parts, on which the
\SourceLeaf{270}
construction of the angles $b',b'',b'''$, etc., depends rationally; then the division of the angle $b'+b''+b'''+\text{etc.}$ into $\beta$ parts; and finally the square root $\surd n$. If we set $\beta=\frac12(n-1)$ at once, those periods evidently coincide with twice the cosines of the angles $\frac{360^0}n$, $2\frac{360^0}n$, $3\frac{360^0}n$, etc., up to $\frac12(n-1)\frac{360^0}n$, so that dividing the circle into $n$ parts depends on dividing the whole circle into $\frac12(n-1)$ parts, dividing an angle that can be constructed after that division into $\frac12(n-1)$ parts, and the radical quantity $\surd n$. If we intend to proceed as far as the sines of the angles $\frac{360^0}n$, etc., one further operation will be needed.
